\documentclass{leveldoc}
% The Foundations: what every level document builds on. Written with the system
% (docs/levels/README.md); level documents cite it as \xeqref{foundations}{eq:...}.

\begin{document}

\tableofcontents

\section{How to read the level documents}\label{sec:reading}

Every level of the game has a document like this one. It derives what happens in the level
from principles found in any textbook of classical electrodynamics, and it carries the
derivation down to numbers: where the particle arrives, when, and how fast. A reader who
knows Maxwell's equations, the Lorentz force and special relativity, and has unlimited time,
should be able to reproduce every number of the game by hand from the level's document, the
documents it cites, and this one.

\emph{By hand} means that nothing is left to a black box. Where the answer has a closed form,
the document derives it. Where it has none (most levels with more than one element), the
document states the exact mathematical problem (equations, initial values, the rule that
ends the flight) and a numerical method, with its error estimate, that a person could carry
out with pencil and paper (§\ref{sec:numerics}). A computer only does that faster.

Each level document has the same nine sections:
\begin{enumerate}[label=\arabic*.]
  \item \emph{What the level asks}: the puzzle, in physical terms, and its goal exactly.
  \item \emph{Setup}: every number of the level file, and the symbols for them.
  \item \emph{Principles and model}: the principles used and every assumption, each with the
    size of what it neglects.
  \item \emph{Derivation}: from the principles to formulas that can be evaluated.
  \item \emph{The reference solution by hand}: the formulas evaluated for the level's
    reference solution, step by step.
  \item \emph{Comparison with the game}: the document's numbers against the game's.
  \item \emph{Other solutions}: what else solves the level, and why.
  \item \emph{What the game computes}: the engine's method for this level, and why it agrees.
  \item \emph{Further reading}.
\end{enumerate}
An appendix reproduces the level file verbatim.

\paragraph{Numbering and references.} Equations carry their document's number: (1.7) is the
seventh equation of Level~1's document, (F.3) the third of this one, in every document.
Sections are cited as §1.4, the engine's own documentation (\textsf{PHYSICS.md}, which
describes the implemented physics and its validation tests) as \textsf{PHYSICS.md}~§5.
A later level cites an earlier one instead of repeating its derivation.

\paragraph{The numbers.} Every number in a document is computed by a short program beside it
(\code{calc.py}) from the level file, with the document's formulas, independently of the
game. The same program compares the results with the game's own (§\ref{sec:verification});
the table in each document's comparison section is its output, so the text and the
comparison cannot drift apart.

\section{Units}\label{sec:units}

The game computes in dimensionless units. A world is fixed by four reference values: a length
$L_0$ (the side of a grid cell), a charge $Q_0$, a mass $M_0$, and the Coulomb constant
$k = 1/(4\pi\varepsilon_0)$, which is set to $1$. The other units follow:
\begin{equation}\label{eq:units}
  T_0 = \sqrt{\frac{M_0 L_0^3}{k Q_0^2}}, \qquad
  E_0 = \frac{k Q_0^2}{L_0}, \qquad
  V_0 = \frac{L_0}{T_0},
\end{equation}
for time, energy and speed. The speed of light becomes a pure number, $c = c_{\mathrm{SI}}/V_0$,
which each level sets (\code{physics.c} in the level file; absent, it is infinite and the
mechanics is Newtonian).

With $k = 1$ the constants of the vacuum are
\begin{equation}\label{eq:vacuum}
  \varepsilon_0 = \frac{1}{4\pi}, \qquad \mu_0 = \frac{4\pi}{c^2}, \qquad
  \text{so that}\quad \varepsilon_0\mu_0 = \frac{1}{c^2}, \quad \frac{\mu_0}{4\pi} = \frac{1}{c^2}.
\end{equation}
These are SI's equations with these values: every formula of an SI textbook holds after
the substitution \eqref{eq:vacuum}.

\begin{remark}
A level with a test charge $q = 10^{-6}$ and sources of $10^{6}$ is not exotic: in the
non-relativistic limit a trajectory depends only on $qQ/m$ and on the energies, so the
puzzle is the same as with charges of order one. The small $q$ keeps the particle's own
radiation negligible (§\ref{sec:radiation}).
\end{remark}

\section{Maxwell's equations and the Lorentz force}\label{sec:maxwell}

\begin{principles}
  \item Maxwell's equations and the Lorentz force law (\Jackson{§§6.1, 6.7};
    \Griffiths{Ch.~7}), in the units of §\ref{sec:units}.
\end{principles}

In these units, Maxwell's equations read
\begin{equation}\label{eq:maxwell}
  \nabla\cdot\vect{E} = 4\pi\rho, \qquad
  \nabla\cdot\vect{B} = 0, \qquad
  \nabla\times\vect{E} = -\pderiv{\vect{B}}{t}, \qquad
  \nabla\times\vect{B} = \frac{4\pi}{c^2}\vect{J} + \frac{1}{c^2}\pderiv{\vect{E}}{t},
\end{equation}
and a charge $q$ moving with velocity $\vect{v}$ feels the force
\begin{equation}\label{eq:lorentz}
  \vect{F} = q\left(\vect{E} + \vect{v}\times\vect{B}\right).
\end{equation}

\subsection{The field of a point charge}\label{sec:coulomb}

For charges at rest the fields do not change, and \eqref{eq:maxwell} gives
$\nabla\times\vect{E} = 0$: there is a potential $\phi$ with $\vect{E} = -\nabla\phi$, and
$\nabla^2\phi = -4\pi\rho$. For a charge $Q$ at the origin, spherically symmetric, the field
is radial, $\vect{E} = E(r)\,\uvec{r}$. Gauss's law (the first equation of
\eqref{eq:maxwell} integrated over a ball of radius $r$ enclosing all of the charge, with the
divergence theorem) gives $4\pi r^2 E(r) = 4\pi Q$:
\begin{equation}\label{eq:coulomb}
  \vect{E}(\vect{x}) = \frac{Q\,\vect{x}}{\abs{\vect{x}}^3}, \qquad
  \phi(\vect{x}) = \frac{Q}{\abs{\vect{x}}},
\end{equation}
the potential chosen to vanish at infinity. The same argument holds for any spherically
symmetric charge of total $Q$ and radius $R$, outside it ($r > R$): the shell theorem. The
equations are linear, so the fields of several charges add:
$\phi(\vect{x}) = \sum_i Q_i/\abs{\vect{x}-\vect{r}_i}$.

\subsection{The force between two charged spheres}\label{sec:spheres}

The game's fixed charges are rigid spheres of radius $R$ (\code{physics.charge\_radius}) with
a spherically symmetric charge, and its particles are rigid spheres of radius $a$ (often $0$)
that do not polarize. Outside each other, the force on the particle is exactly
\begin{equation}\label{eq:sphereforce}
  \vect{F} = q\,\vect{E}_Q(\vect{x}) = \frac{qQ\,(\vect{x}-\vect{r})}{\abs{\vect{x}-\vect{r}}^3},
\end{equation}
with $\vect{x}$ and $\vect{r}$ the centres: by the shell theorem the particle's field outside
it is that of a point charge at its centre, so the force it exerts on the source is that of a
point charge; by Newton's third law for the pairwise central Coulomb forces between their
elements, the force on the particle is the opposite, which is \eqref{eq:sphereforce}.

\subsection{Static magnetic fields}

Steady currents make a magnetic field by the Biot--Savart law, in these units
\begin{equation}\label{eq:biotsavart}
  \vect{B}(\vect{x}) = \frac{1}{c^2}\oint \frac{I\,\dd\vect{l}\times(\vect{x}-\vect{l})}{\abs{\vect{x}-\vect{l}}^3}
\end{equation}
(\Jackson{§5.2}, with $\mu_0/4\pi = 1/c^2$). The levels with magnets and coils derive what
they need from it.

\section{The motion of a charge}\label{sec:mechanics}

\begin{principles}
  \item Special relativity: the momentum of a particle of mass $m$ is
    $\vect{p} = \gamma m\vect{v}$, its energy $\gamma mc^2$ (\LandauII{§9}).
  \item Newton's second law in the form $\dd\vect{p}/\dd t = \vect{F}$, which holds in
    relativity with the Lorentz force (\LandauII{§17}; \Jackson{§12.1}).
\end{principles}

The state of a particle is its position $\vect{x}$ and momentum $\vect{p}$, and
\begin{equation}\label{eq:eom}
  \deriv{\vect{x}}{t} = \vect{v} = \frac{\vect{p}}{\gamma m}, \qquad
  \deriv{\vect{p}}{t} = q\left(\vect{E}(\vect{x},t) + \vect{v}\times\vect{B}(\vect{x},t)\right), \qquad
  \gamma = \sqrt{1 + \frac{\abs{\vect{p}}^2}{m^2c^2}}.
\end{equation}
Then $\abs{\vect{v}} = \abs{\vect{p}}c/\sqrt{m^2c^2+\abs{\vect{p}}^2} < c$ for every finite
momentum. As $c\to\infty$, $\gamma\to 1$ and \eqref{eq:eom} becomes Newton's mechanics.

\paragraph{The launch.} A particle launched with kinetic energy
$T_0 = (\gamma_0-1)mc^2$ has
\begin{equation}\label{eq:launch}
  \gamma_0 = 1 + \frac{T_0}{mc^2}, \qquad
  \abs{\vect{p}_0} = mc\sqrt{\gamma_0^2-1} = \sqrt{\frac{T_0^2}{c^2} + 2mT_0},
\end{equation}
along the launch direction.

\subsection{Energy}\label{sec:energy}

The energy $\mathcal{E} = \gamma mc^2 = \sqrt{m^2c^4 + \abs{\vect{p}}^2c^2}$ changes at the
rate $\dd\mathcal{E}/\dd t = c^2\,\vect{p}\cdot\dot{\vect{p}}/\mathcal{E} = \vect{v}\cdot\dot{\vect{p}}$.
The magnetic force does no work ($\vect{v}\cdot(\vect{v}\times\vect{B}) = 0$). In a static
electric field $\vect{E} = -\nabla\phi$ the rest is
$q\,\vect{v}\cdot\vect{E} = -q\,\vect{v}\cdot\nabla\phi = -q\,\dd\phi(\vect{x}(t))/\dd t$, so
\begin{equation}\label{eq:energy}
  W = \gamma mc^2 + q\,\phi(\vect{x})
\end{equation}
is constant along the flight. Equivalently, the kinetic energy $T = (\gamma-1)mc^2$ satisfies
$T + q\phi(\vect{x}) = T_0 + q\phi(\vect{x}_0)$.

\subsection{Angular momentum}\label{sec:angmom}

The angular momentum about a point $\vect{r}$ is
$\vect{L} = (\vect{x}-\vect{r})\times\vect{p}$, and
$\dd\vect{L}/\dd t = \vect{v}\times\vect{p} + (\vect{x}-\vect{r})\times\vect{F}$. The first term
vanishes ($\vect{p}\parallel\vect{v}$). If the force points along $\vect{x}-\vect{r}$, as the
field of a single charge at $\vect{r}$ does \eqref{eq:sphereforce}, the second vanishes too:
\begin{equation}\label{eq:angmom}
  \vect{L} = (\vect{x}-\vect{r})\times\vect{p} \quad\text{is constant in a central field,}
\end{equation}
in relativity as in Newton's mechanics.

\subsection{Motion in the plane}\label{sec:plane}

Most levels lie in the plane $z = 0$: their charges sit in it and their particles are launched
along it. The field of charges in the plane is mirror-symmetric, $E_z(x, y, 0) = 0$, and a
magnetic field of currents in the plane is perpendicular to it there, so
$(\vect{v}\times\vect{B})_z = 0$ for $\vect{v}$ in the plane. Then $z(t) = 0$, $p_z(t) = 0$
solves \eqref{eq:eom}, and by the uniqueness of solutions of ordinary differential equations
with smooth right-hand sides, it is the solution. Plane polar coordinates $(r,\varphi)$ about
a point serve wherever a level has one centre.

\section{The board}\label{sec:board}

\paragraph{The grid.} Elements sit on the nodes of a grid: node $(i,j,k)$ is the point
$(i,j,k)/s$, with $s$ the grid's subdivision (\code{grid.subdivision}, usually 1). A planar
level (\code{grid.nz} = 0) has nodes $0\le i\le n_x$, $0\le j\le n_y$ in the plane $z = 0$.

\paragraph{The world.} The world is the grid's box enlarged by one cell on every side:
$-1\le x\le n_x/s+1$, $-1\le y\le n_y/s+1$ and $-1\le z\le n_z/s+1$. A flight that leaves it
ends, lost.

\paragraph{Shots.} A shot launches a particle (charge $q$, mass $m$, radius $a$) from a node,
along a direction (normalised), with kinetic energy $T_0$, at $t = 0$.

\paragraph{Fixed charges.} A charge element at node $\vect{r}$ is a sphere of radius $R$
(\code{physics.charge\_radius}) with charge $Q$. A flight that comes within $R + a$ of
$\vect{r}$ ends, a collision.

\paragraph{The detector.} Each shot has a detector, an axis-aligned box (in a planar level a
slab through the plane, $-1\le z\le 1$: in the plane, a rectangle). The first entry decides:
if the detector sets conditions on the arriving particle (a cone of directions, a band of
kinetic energies), the state at the moment of entry must meet them.

\paragraph{The time limit.} A flight that has not ended by \code{physics.t\_max} times out.

\paragraph{The outcome.} The earliest of these events ends the flight: the particle touches
a charge or another obstacle, leaves the world, enters its detector, or runs out of time.
Should two happen at the same instant, an obstacle counts before the world's edge, and that
before the detector. A level is solved when every shot's flight arrives in its detector and
meets its conditions (with disturbances, in every one of them), and the game has verified it
(§\ref{sec:verification}).

The levels with gates, beams of many particles, conductors, coils, waves, circuits or the
tubes' $z$-invariant world add their rules in the first level of their arc; its document
derives them, and the later ones cite it.

\section{What the game verifies}\label{sec:verification}

The game integrates every flight twice, at two error tolerances of its integrator (the
\emph{preview}, usually $10^{-10}$, and the \emph{verification}, $10^{-12}$:
\code{physics.tolerances}), and computes for every boundary the closest approach to it, or
how deep the flight went in. A result is \emph{verified} when both runs agree on the
outcome and every margin is at least ten times the two runs' difference
(\physref{7}). The difference estimates the preview's error; the verification's error is
about a hundred times smaller. Measured on orbits with exact solutions, the verified run's
positions are accurate to about $10^{-10}$ cells (\physref{5}).

So when a level's physics has an exact solution, the document's numbers and the game's agree
to about $10^{-10}$ (in cells and time units); a tolerance of $10^{-9}$ in the comparison is
then a test of both. Where the document approximates, its tolerance is its own error
estimate.

\section{Radiation and other neglected effects}\label{sec:radiation}

An accelerated charge radiates. A particle of charge $q$ with velocity $\vect{v}$ and
acceleration $\vect{a}$ radiates the power (Liénard's generalization of Larmor's formula,
\Jackson{§14.2}, in the units of §\ref{sec:units})
\begin{equation}\label{eq:larmor}
  P = \frac{2}{3}\,\frac{q^2}{c^3}\,\gamma^6\left(\abs{\vect{a}}^2 - \frac{\abs{\vect{v}\times\vect{a}}^2}{c^2}\right).
\end{equation}
Unless a level includes radiation reaction (its document then derives it), the game ignores
the energy this carries away; it integrates $\int P\,\dd t$ along each flight as a check, and
every shipped level keeps it below $10^{-10}\,T_0$. By hand: near a charge $Q$ at distance $r$
the acceleration is about $qQ/(m r^2)$ for a duration of about $r/v$, so the radiated
fraction of the kinetic energy is of order $(q^2/mc^2)/r$ times powers of $v/c$ and $qQ/(rT)$.
With $q = 10^{-6}$, $m = 1$ this is far below $10^{-10}$.

Other effects the game leaves out, unless a level says otherwise: the recoil of the fixed
charges (held in place), the polarization of the particle, and the particle's field acting
back on the sources (\physref{8}).

\section{Computing by hand}\label{sec:numerics}

These methods carry out what a closed form cannot. Each comes with an error estimate, so that
a hand computation can be made as accurate as a comparison needs
(\HairerI{§§II.1, II.4} for the Runge--Kutta methods).

\paragraph{Roots.} To solve $f(s) = 0$ on an interval $[s_a, s_b]$ where $f$ changes sign:
\emph{bisection} halves the interval, keeping the half where the sign changes; after $n$
halvings the root is known to $(s_b-s_a)/2^n$. \emph{Newton's method},
\begin{equation}\label{eq:newton}
  s_{k+1} = s_k - \frac{f(s_k)}{f'(s_k)},
\end{equation}
converges quadratically near a simple root (the number of correct digits roughly doubles each
step); stop when $\abs{s_{k+1}-s_k}$ is below the accuracy needed, and fall back to bisection
if a step leaves the interval.

\paragraph{Integrals.} The composite Simpson rule with an even number $n$ of steps
$h = (s_b-s_a)/n$,
\begin{equation}\label{eq:simpson}
  \int_{s_a}^{s_b} f\,\dd s \approx \frac{h}{3}\left(f_0 + 4f_1 + 2f_2 + 4f_3 + \dots + 4f_{n-1} + f_n\right),
\end{equation}
has an error $-(s_b-s_a)h^4 f^{(4)}(\xi)/180$; halving $h$ and comparing estimates the error as
$(S_h - S_{2h})/15$.

\paragraph{Ordinary differential equations.} For $\dd y/\dd t = f(t,y)$, the classical
Runge--Kutta method advances $y_n$ by a step $h$:
\begin{equation}\label{eq:rk4}
  \begin{aligned}
    k_1 &= f(t_n, y_n), &
    k_2 &= f\!\left(t_n + \tfrac{h}{2}, y_n + \tfrac{h}{2}k_1\right), \\
    k_3 &= f\!\left(t_n + \tfrac{h}{2}, y_n + \tfrac{h}{2}k_2\right), &
    k_4 &= f(t_n + h, y_n + hk_3), \\
    y_{n+1} &= y_n + \tfrac{h}{6}\left(k_1 + 2k_2 + 2k_3 + k_4\right). &&
  \end{aligned}
\end{equation}
Its global error is of order $h^4$: integrating with $h$ and with $h/2$ and comparing
estimates the error of the finer run as $(y_{h/2}-y_h)/15$. For a flight, $y = (\vect{x},\vect{p})$
and $f$ is the right-hand side of \eqref{eq:eom}.

\paragraph{Where a flight ends.} An event (entering a box, touching a sphere) is the first
zero of a function $g(\vect{x})$ of the position: the signed distance to the boundary,
positive outside. After each step, check the sign of $g$. In the step where it changes, the
cubic polynomial that matches the position and velocity at both ends of the step (cubic
Hermite interpolation; its error is of order $h^4$, like the method's) gives $\vect{x}(t)$ in
between, and bisection or Newton's method on $g(\vect{x}(t))$ finds the instant. A step longer
than the boundary's distance divided by the speed can jump over a thin corner; keep steps
shorter than that near boundaries.

\paragraph{What the game does instead.} The engine integrates with an eighth-order
Runge--Kutta method with step-size control (Dormand--Prince 8(5,3), \HairerI{§II.10}),
locates events on its seventh-order interpolant with a search that cannot miss a crossing,
and verifies every result as in §\ref{sec:verification} (\physref{5}, \physref{6}). For
a hand computation the simpler methods above suffice: only their error estimates matter.

\section{Notation}\label{sec:notation}

\begin{center}
\small
\begin{tabular}{@{}ll@{}}
\toprule
Symbol & Meaning \\
\midrule
$q$, $m$, $a$ & the particle's charge, mass and radius \\
$Q$, $R$ & a fixed charge and its radius (\code{physics.charge\_radius}) \\
$\kappa = qQ$ & the coupling of particle and charge ($k = 1$); $\kappa > 0$ repels \\
$c$ & the speed of light (\code{physics.c}) \\
$T$, $T_0$ & kinetic energy, at launch \\
$\gamma$, $\vect{p}$, $\vect{v}$ & Lorentz factor, momentum, velocity \\
$W$ & the conserved energy \eqref{eq:energy} \\
$\vect{L}$, $L$ & angular momentum \eqref{eq:angmom}, its magnitude \\
$r$, $\varphi$ & plane polar coordinates about a charge \\
$u = 1/r$ & the inverse radius (orbits) \\
$t_\mathrm{e}$, $\vect{x}_\mathrm{e}$, $\vect{p}_\mathrm{e}$ & time, position and momentum at the event that ends a flight \\
\bottomrule
\end{tabular}
\end{center}

\section{Further reading}\label{sec:reading-list}

\begin{itemize}
  \item \Jackson{Chs.~1, 5, 6, 12 and 14}: electrostatics, magnetostatics, Maxwell's
    equations, relativistic dynamics of charges, radiation of moving charges.
  \item \LandauII{§§8--17 and 39}: relativistic mechanics, charges in fields, motion in a
    Coulomb field.
  \item \Griffiths{Chs.~2, 5, 7, 12}: the same at an undergraduate level, in SI units.
  \item \HairerI{Ch.~II}: Runge--Kutta methods, error estimation and step-size control.
\end{itemize}

\end{document}
